\section{Introduction}\label{sec:offline-intro}

An offline payment is one made while both the payer and the payee are
disconnected from the network, and settled only once the token that changed
hands is eventually deposited. Offline operation matters for availability and
for accessibility. A payment instrument that stops working during a blackout, or
in a place with poor network coverage, is not a replacement for cash.

We work in the setting of central bank digital currency (CBDC). A CBDC is a form
of central bank money on par with cash and central bank reserves, represented
digitally as tokens. Each token encodes an owner and a value.  Retail CBDC is
the variant available to the general public rather than to financial
institutions. It is the most demanding to realise. It must be resilient, scale
to a national population, protect the privacy of its users, and support payments
made offline~\cite{cbdctracker}.

The difficulty is that the payee cannot check, at the time of the payment,
whether the token they are receiving has already been spent. Without that check,
double spending is trivial. Responses divide into prevention and deterrence. A
wallet application can delete a token as soon as it is sent, but an attacker who
tampers with the application bypasses the restriction. A secure element enforces
the same rule in hardware, and is much harder to subvert. Deterrence instead
identifies the double spender once the conflict surfaces at deposit, and relies
on the consequences of being identified.

We combine the two, because neither is sufficient alone. A secure element offers
no recourse once it is broken, and identification after the fact does not stop
an attacker from spending a token many times before any deposit occurs.
Together, the secure element makes double spending expensive and impractical to
mount at scale, and identification ensures that an attack that does succeed can
be attributed to the party responsible and settled against them.

An offline payment system has operational roles that already exist. A central
bank issues and redeems tokens and controls the money supply. Intermediaries
hold users' accounts, onboard them, and route their requests. A registration
authority enrols users following established know your customer processes. An
auditor is empowered to attribute a fraudulent payment. Each of these is
necessary for the system to function. None of them needs to learn who paid whom.
We keep these roles and constrain what they learn. Payments are anonymous to
everyone who sees a transcript. The central bank and the intermediaries cannot
link a withdrawal to a deposit. The auditor recovers an identity only from a
payment that double spends a token.

We give a transferable offline payment scheme for this setting. Users obtain a
long term credential once, then pay one another offline by extending a chain of
anonymous signatures over a token. Each payment carries a verifiable encryption
of the payer's identifier under the auditor's key. The cryptography this
requires is too heavy for an off the shelf secure element, so we split each
payment in two. The secure element holds the signing key and enforces that a
token is used once. The user's mobile device performs the rest.

\subsection{Related work}\label{sec:offline-related}

\paragraph{Original ecash.} Ecash has been studied for several decades,
beginning with Chaum's solutions, which focus on the privacy of users and on
double spending detection~\cite{C:Chaum82,C:ChaFiaNao88}. Privacy is achieved
through RSA blind signatures, which break the link between withdrawals and
deposits, while double spending detection is realised using unique serial
numbers. Although efficient, these solutions guarantee neither transferability
nor theft prevention.

\paragraph{Divisible ecash.} Divisible ecash allows a user to withdraw a token
and split it among different recipients across multiple independent
transactions~\cite{C:OkaOht91,C:Okamoto95}. This comes at the cost of higher
complexity for withdrawals, payments, and deposits, which scale with the token
value. This is partially addressed by work that achieves constant computation
and communication complexity for withdrawal and payment, though deposits still
carry large communication and computational costs~\cite{PKC:CPST15,ACNS:CPST15}.
While divisibility is a useful feature, we leave it as future work and focus
instead on transferability.

\paragraph{Transferable ecash.} Some schemes present ecash that is both
transferable and divisible~\cite{C:OkaOht91,EC:ChaPed92}, at the cost of high
computational complexity. We instead focus on efficiency and scalability.
Another system maintains a constant token size regardless of the number of
transfers, but requires user interaction to identify double
spenders~\cite{CANS:FucPoiVer09}. When double spending is detected, all previous
owners of the token also become identifiable. There is also a scheme where token
size grows linearly with the number of transfers, but which reveals only the
double spender's identity~\cite{PKC:BCFK15}. Owing to the cost of this approach,
we instead rely on a trusted auditor, which can be distributed, to deanonymise
payments as needed.

Some schemes introduce coin transparency, which ensures a payee cannot tell
whether they are receiving a token they previously
owned~\cite{ACNS:CanGou08,PKC:BauFucQia21}. One such scheme achieves this
through extensive use of zero knowledge proofs together with a trusted party
that tracks every token in the system and can identify double spending
attempts~\cite{AFRICACRYPT:BCFGST11}. While this captures the strongest level of
privacy, we do not target coin transparency, owing to its computational overhead
and to the fact that physical cash does not satisfy it either.

The bounded anonymity model keeps transactions anonymous as long as each user
interacts with a given merchant no more than a fixed number of times. Beyond
this, the user becomes identifiable~\cite{SCN:CamHohLys06}. We instead let the
conditions under which anonymity is revoked be defined outside of the protocol
itself.

\paragraph{Offchain settlement.} Payment channels such as the Lightning network
and Teechain provide offchain payments rather than offline
payments~\cite{poon2016bitcoin,lind2019teechain}. The Lightning network was the
first payment channel network designed for Bitcoin~\cite{poon2016bitcoin}.
Users commit funds to an account controlled by signatures from both payer and
payee, then exchange small transactions offchain. Only the final batched
transaction is committed to the ledger, which improves latency and scalability.
A network of such channels lets users route transactions through counterparties
with whom they do not share an account directly. Onion routing offers some
notion of privacy, though it has been shown that ordinary use of the network can
bypass these privacy aims~\cite{FC:KYPKDM21}. Teechain instead has users fund an
account controlled by a committee of treasurers, who then use that value
offchain~\cite{lind2019teechain}. The treasurers store their share of the
private key in trusted hardware, which raises the cost of stealing funds.
Teechain does not target privacy, and is motivated by reducing the onchain
storage costs of every transaction rather than by enabling offline payments.

\paragraph{CBDC offline payments.} Several reports motivate the need for CBDC
offline payments and describe current technologies. Project
Polaris\footnote{Project Polaris: secure and resilient CBDC systems, offline and
online.  \url{https://www.bis.org/about/bisih/topics/cbdc/polaris.htm}.} from
the Bank for International Settlements sets out design goals for offline payment
systems, distinguishing three types. Fully offline systems settle offline.
Intermittently offline systems settle offline in general, but require occasional
synchronisation with the ledger, which might be triggered once the total value
transferred crosses a risk threshold. Staged offline systems exchange value
offline but do not settle it until the payee comes online. Our scheme is fully
offline in this sense: a payee accepts a payment on the strength of the
transcript alone, and needs no contact with any other party either to accept it
or to pay it onward.

The deployed and announced systems we are aware of are all staged offline.
Sweden has piloted its e-krona project on the Corda platform, where locally
stored e-krona can be used offline and, like cash, is not recoverable if the
wallet is lost, but transactions formed offline are not settled until the device
goes online, as double spending is checked
online~\cite{brown2016corda,soderberg2019krona}. The People's Bank of China has
designed a system for digital payments in which transaction values are hidden;
no technical specification has been published, and press reporting suggests that
user anonymity and transaction unlinkability are not implemented and that
offline payments are not specifically addressed.\footnote{Reported at
\url{https://www.nytimes.com/2021/03/01/technology/china-national-digital-currency.html}.}
The Bank of Ghana's design document for the digital Cedi (eCedi) highlights
offline payments as a key focus area while emphasising privacy alongside
regulatory compliance.\footnote{Reported at
\url{https://www.bog.gov.gh/news/design-paper-of-the-digital-cedi-ecedi/}.} The
US Federal Reserve Board's report distinguishes hybrid payment systems, where
the payee must come online to settle, from fully offline systems, where this is
not required, and finds no production examples of the
latter~\cite{aboulaiz2024offline}. It highlights the digital Rupee from the
National Payments Corporation of India, which allows offline transactions from a
smartphone app, settled once the app reconnects,\footnote{Reported at
\url{https://rbi.org.in/scripts/FAQView.aspx?Id=169}.} and Pix in Brazil, which
allows the payer to be offline when creating a transaction, but requires the
payee's point of sale terminal to be online to authorise it~\cite{lobo2021pix}.

\paragraph{Decentralised payments.} Peredi, Platypus, and UTT give privacy
preserving token schemes that distribute trust via threshold cryptography while
providing accountability and compliance guarantees, but do not address the
offline setting~\cite{CCS:KiaKohSar22,CCS:WKDC22,EPRINT:TBAAGP22}. OPERA enables
offline payments and prevents double spending using one-time readable memory,
but this technology is not widely available in practice~\cite{park2017opera}.
Rial and Piotrowska distribute the role of the bank via threshold issuance and
consider divisible offline payments, while we focus on efficient
transferability~\cite{PoPETS:RiaPio23}. PayOff offers a transferable offline
CBDC solution using more complex cryptography than ours~\cite{Beer2024}. It
assigns the central bank responsibility for covering double spent values, which
can inflate the total value in the system. We instead equip users with a secure
element, to prevent double spending rather than only detect it, and settle any
double spent value against the party the auditor identifies.

\Cref{tab:offline-related-comparison} compares our scheme with the offline
payment schemes closest to it. All of the ecash schemes listed are offline in
the sense above, so the axes that distinguish them are transferability,
divisibility, what happens once a conflict surfaces, and what the scheme asks of
the user's wallet. We target transferability rather than divisibility. We
identify the double spender rather than absorbing the loss or identifying every
previous holder of the token. We assume a secure element, so that prevention and
identification reinforce one another.

\begin{table}[t] \centering \small
%
  \begin{tabular}{l|c|c|c|c}
%
    \toprule Scheme & Transferable & Divisible & Double spending & Hardware \\
\midrule \hline \cite{C:Chaum82,C:ChaFiaNao88} & No & No & Detected & None \\
%
    \cite{C:OkaOht91,C:Okamoto95} & No & Yes & Detected & None \\
%
    \cite{PKC:CPST15,ACNS:CPST15} & No & Yes & Detected & None \\
%
    \cite{CANS:FucPoiVer09} & Yes & No & All holders & None \\
%
    \cite{PKC:BCFK15} & Yes & No & Spender only & None \\
%
    \cite{park2017opera} & -- & -- & Prevented & OTR memory \\
%
    \cite{PoPETS:RiaPio23} & No & Yes & Detected & None \\
%
    \cite{Beer2024} & Yes & No & Absorbed & None \\
%
    Our scheme & Yes & No & Spender only & Secure element \\ \bottomrule
%
  \end{tabular}
%
  \caption{Comparison with existing offline payment schemes. The double spending
column records what happens once a conflict is detected at deposit. Detected
means the conflict is observed but no party is named, all holders means every
previous owner of the token is identified, spender only means the double spender
alone is identified, absorbed means the issuer covers the double spent value,
and prevented means double spending is stopped rather than resolved after the
fact. The hardware column records the hardware assumption the scheme places on
the user's wallet, where OTR denotes one-time readable memory. Entries marked --
are not addressed by the cited work.} \label{tab:offline-related-comparison}
\end{table}

\paragraph{Our contribution.} We give a transferable, privacy preserving offline
payment scheme. The contribution is not a new primitive but a combination, and
consists of the following.

\begin{itemize}

  \item A syntax for offline payment schemes and a simulation based security
definition for them, given as a single ideal functionality that captures
unforgeability, theft prevention, double spending identification,
non-frameability, anonymity, unlinkability, and conservation of value at once.

  \item A construction that splits the cost of authorising a payment between a
lightweight secure element and a more powerful mobile device, so that the secure
element need only perform operations comparable in cost to producing an ordinary
digital signature. Anonymous credentials, blind signatures, and zero knowledge
proofs of knowledge hide the identities of payers and payees from the central
bank and from the intermediaries, and a verifiable encryption under an auditor's
key allows double spenders to be identified without compromising the privacy of
honest users.

  \item An instantiation from BBS+ and Pointcheval-Sanders signatures over
pairing friendly groups, with a proof that it securely realises the ideal
functionality.

  \item An implementation on commodity hardware, integrated into a mobile wallet
on iOS and Android and a Javacard applet on a smart card, and benchmarked end to
end.

\end{itemize}

\paragraph{Layout.} \Cref{sec:offline-definitions} describes the participants
and trust assumptions, gives the syntax of an offline payment scheme, and
formalises the security goals as an ideal functionality.
\Cref{sec:offline-construction} gives the construction, first informally and
then in full. \Cref{sec:offline-instantiation} instantiates the building blocks
it relies on. \Cref{sec:offline-security} proves both the construction and its
instantiation secure. \Cref{sec:offline-evaluation} evaluates our
implementation, and \Cref{sec:offline-summary} concludes.
